\documentclass[a4paper,11pt]{article}

\usepackage[utf8]{inputenc}
\usepackage[T1]{fontenc}
\usepackage{lmodern}
\usepackage[ngerman]{babel}
\usepackage{amsmath,amssymb}
\usepackage[a4paper,top=2.4cm,bottom=2cm,left=2cm,right=2cm]{geometry}
\usepackage{xcolor}
\usepackage{tikz}
\usepackage{enumitem}
\usepackage{fancyhdr}

% ---------- Dark Mode ----------
\definecolor{bgdark}{HTML}{1E1E1E}
\definecolor{fgtext}{HTML}{E6E6E6}
\definecolor{gridcol}{HTML}{4A4A4A}
\definecolor{accent}{HTML}{6FB7FF}
\definecolor{easy}{HTML}{7FD28A}
\definecolor{med}{HTML}{F0C45A}
\definecolor{hard}{HTML}{F08A7A}
\pagecolor{bgdark}
\color{fgtext}

% ---------- Kopf-/Fußzeile ----------
\pagestyle{fancy}
\fancyhf{}
\renewcommand{\headrulewidth}{0pt}
\renewcommand{\footrulewidth}{0pt}
\fancyfoot[C]{\textcolor{gridcol}{\small BWL -- Übungsaufgaben 01 \;|\; Seite \thepage}}

% ---------- Karo-Ausfüllfeld ----------
\newcommand{\karo}[1]{\par\vspace{2mm}\noindent%
  \begin{tikzpicture}
    \draw[step=5mm, gridcol, very thin] (0,0) grid (\linewidth,#1);
  \end{tikzpicture}\par}

% ---------- Schwierigkeits-Badge ----------
\newcommand{\diff}[2]{\colorbox{#1}{\textcolor{bgdark}{\small\textbf{#2}}}}

% ---------- Aufgaben-Kopf (immer neue Seite) ----------
\newcommand{\aufgabe}[3]{\clearpage%
  \noindent{\color{accent}\Large\bfseries #1}\hfill #2\par
  \vspace{1mm}\noindent\textcolor{gridcol}{\rule{\linewidth}{0.4pt}}\par
  \vspace{1mm}\noindent#3\par}

\setlength{\parindent}{0pt}
\linespread{1.05}

\begin{document}

% ===================== TITELSEITE =====================
\thispagestyle{empty}
\vspace*{1cm}
\begin{center}
  {\color{accent}\Huge\bfseries BWL -- Übungsaufgaben}\\[4pt]
  {\Large Ausgewählte Themen der Digitalen Betriebswirtschaftslehre}\\[10pt]
  {\large Klausurvorbereitung -- Aufgabenblatt 01}\\[2pt]
  {\small (Aufgaben zum Ausfüllen -- Lösungen separat)}
\end{center}

\vspace{1.4cm}
\noindent\textcolor{gridcol}{\rule{\linewidth}{0.4pt}}\\[3mm]
{\large Name:\enspace\textcolor{gridcol}{\rule{6cm}{0.4pt}}\hfill Datum:\enspace\textcolor{gridcol}{\rule{4cm}{0.4pt}}}\\[3mm]
\noindent\textcolor{gridcol}{\rule{\linewidth}{0.4pt}}

\vspace{1.2cm}
\noindent{\color{accent}\large\bfseries Hinweise}
\begin{itemize}[leftmargin=*]
  \item 12 Aufgaben, gemischte Schwierigkeit:\enspace
        \diff{easy}{leicht}\enspace \diff{med}{mittel}\enspace \diff{hard}{schwer}.
  \item Es zählt \textbf{Verständnis statt Auswendiglernen} -- begründe deine Antworten.
  \item Bei Rechenaufgaben: Rechenweg angeben, Ergebnisse interpretieren.
  \item Benötigte Formeln sind in der jeweiligen Aufgabe angegeben.
  \item Jede Aufgabe beginnt auf einer neuen Seite, darunter ist Platz zum Ausfüllen.
\end{itemize}

\vspace{0.6cm}
\noindent{\color{accent}\large\bfseries Abgedeckte Themen}\\[1mm]
VL1 Einführung \;$\cdot$\; VL2 Märkte \;$\cdot$\; VL4 Strategie \;$\cdot$\;
VL5 Marketing \;$\cdot$\; VL7 Beschaffung/SCM \;$\cdot$\; VL8 Personal

% ===================== AUFGABE 1 =====================
\aufgabe{Aufgabe 1 -- Ökonomisches Prinzip}{\diff{easy}{leicht} \;(8 P)}{%
VL1 -- Einführung}

\begin{enumerate}[label=\textbf{\alph*)}, leftmargin=*]
  \item Ordnen Sie jede Situation dem \textbf{Minimum-}, \textbf{Maximum-} oder
        \textbf{Optimum-Prinzip} zu und begründen Sie kurz:
  \begin{enumerate}[label=(\arabic*), leftmargin=*]
    \item Eine Spedition will mit einem fest vorgegebenen Budget von 50.000\,€ möglichst
          viele Pakete ausliefern.
    \item Ein Cloud-Anbieter will eine vertraglich fest zugesagte Rechenleistung mit
          möglichst geringem Energieeinsatz bereitstellen.
    \item Ein Start-up sucht das günstigste Verhältnis aus Entwicklungskosten und
          Funktionsumfang einer App.
  \end{enumerate}
  \item Erläutern Sie den Unterschied zwischen \textbf{Effizienz} und \textbf{Effektivität}.
  \item Ein Team programmiert ein Feature mit minimalem Ressourceneinsatz -- das am Ende
        aber kein Kunde nutzen will. Ist das Vorgehen effizient, effektiv, beides oder
        keines? Begründen Sie.
\end{enumerate}
\karo{13cm}

% ===================== AUFGABE 2 =====================
\aufgabe{Aufgabe 2 -- Produktivität, Wirtschaftlichkeit, Rentabilität}{\diff{med}{mittel} \;(12 P)}{%
VL1 -- Einführung}

Ein Hersteller von Sensoren fertigt in einem Jahr \textbf{20.000 Sensoren} und setzt dafür
\textbf{4.000 Arbeitsstunden} ein. Der gesamte \textbf{Aufwand} beträgt \textbf{480.000\,€},
der \textbf{Ertrag} (Umsatzerlös) \textbf{600.000\,€}. Das durchschnittlich eingesetzte
\textbf{Kapital} beträgt \textbf{400.000\,€}.

\smallskip
\textit{Formeln:}\;
$\text{Produktivität}=\dfrac{\text{Output}}{\text{Input}}$,\quad
$\text{Wirtschaftlichkeit}=\dfrac{\text{Ertrag}}{\text{Aufwand}}$,\quad
$\text{Gewinn}=\text{Ertrag}-\text{Aufwand}$,\quad
$\text{Rentabilität}=\dfrac{\text{Gewinn}}{\text{Kapital}}\cdot 100$.

\begin{enumerate}[label=\textbf{\alph*)}, leftmargin=*]
  \item Berechnen Sie die \textbf{Arbeitsproduktivität}, die \textbf{Wirtschaftlichkeit},
        den \textbf{Gewinn} und die \textbf{Rentabilität}.
  \item Interpretieren Sie die Wirtschaftlichkeit: Arbeitet das Unternehmen wirtschaftlich?
  \item Die Produktivität soll um \textbf{20\,\%} steigen. Bestimmen Sie den Zielwert und
        nennen Sie zwei Maßnahmen -- ordnen Sie diese dem Minimum- bzw.\ Maximum-Prinzip zu.
\end{enumerate}
\karo{12cm}

% ===================== AUFGABE 3 =====================
\aufgabe{Aufgabe 3 -- SMART-Ziele und Zielbeziehungen}{\diff{easy}{leicht} \;(8 P)}{%
VL1 -- Einführung}

\begin{enumerate}[label=\textbf{\alph*)}, leftmargin=*]
  \item Ein Unternehmen formuliert: \glqq Wir wollen bald deutlich besser im Markt sein.\grqq{}
        Erklären Sie, warum dieses Ziel \textbf{nicht SMART} ist, und formulieren Sie es als
        \textbf{SMART}-Ziel neu (alle fünf Kriterien sichtbar).
  \item Beurteilen Sie die folgenden Zielpaare als \textbf{komplementär},
        \textbf{konkurrierend} oder \textbf{indifferent}:
  \begin{enumerate}[label=(\arabic*), leftmargin=*]
    \item Senkung der Lieferzeit \;$\leftrightarrow$\; Erhöhung der Kundenzufriedenheit
    \item Maximierung des kurzfristigen Gewinns \;$\leftrightarrow$\; hohe Investitionen in Forschung
    \item Umstellung der Kantine auf vegane Gerichte \;$\leftrightarrow$\; Erhöhung der Server-Verfügbarkeit
  \end{enumerate}
\end{enumerate}
\karo{14cm}

% ===================== AUFGABE 4 =====================
\aufgabe{Aufgabe 4 -- Stakeholder und Principal-Agent}{\diff{med}{mittel} \;(11 P)}{%
VL1 -- Einführung}

Die \textit{DataFlow GmbH} entwickelt eine SaaS-Plattform. Das Management wird von den
Eigentümern eingesetzt und geführt.

\begin{enumerate}[label=\textbf{\alph*)}, leftmargin=*]
  \item Nennen Sie \textbf{vier Stakeholder} (zwei interne, zwei externe) und jeweils deren
        \textbf{Anspruch} gegenüber dem Unternehmen.
  \item Erklären Sie den Unterschied zwischen \textbf{Stakeholder-} und \textbf{Shareholder-Ansatz}.
  \item Wenden Sie die \textbf{Principal-Agent-Theorie} an: Wer ist Principal, wer Agent?
        Worin besteht das Grundproblem (Stichwort \textbf{Informationsasymmetrie})?
        Nennen Sie \textbf{zwei Mechanismen}, um es zu verringern.
\end{enumerate}
\karo{12.5cm}

% ===================== AUFGABE 5 =====================
\aufgabe{Aufgabe 5 -- Marktgleichgewicht}{\diff{med}{mittel} \;(12 P)}{%
VL2 -- Unternehmen und Märkte}

Auf einem Wettbewerbsmarkt für USB-Ladegeräte gelten:
\[
\text{Nachfrage:}\quad Q_D = 1800 - 30\,P
\qquad\qquad
\text{Angebot:}\quad Q_S = 200 + 50\,P
\]
($Q$ = Menge, $P$ = Preis in €).

\begin{enumerate}[label=\textbf{\alph*)}, leftmargin=*]
  \item Bestimmen Sie \textbf{Gleichgewichtspreis} und \textbf{Gleichgewichtsmenge}.
  \item Der Staat setzt einen \textbf{Höchstpreis von $P=10$\,€} fest. Berechnen Sie Angebot
        und Nachfrage. Liegt ein Angebots- oder Nachfrageüberschuss vor (Höhe)? Was passiert
        am Markt?
  \item Ein Trend verschiebt die Nachfrage zu $Q_D = 2200 - 30\,P$. Bestimmen Sie das
        \textbf{neue Gleichgewicht} und erklären Sie die Veränderung.
  \item Welche Rolle spielen \textbf{Daten} (z.\,B.\ Verkaufszahlen, Suchanfragen) für die
        Schätzung der Nachfrage -- und warum ist das ein Wettbewerbsvorteil?
\end{enumerate}
\karo{11cm}

% ===================== AUFGABE 6 =====================
\aufgabe{Aufgabe 6 -- Marktformen und digitale Plattformen}{\diff{easy}{leicht} \;(9 P)}{%
VL2 -- Unternehmen und Märkte}

\begin{enumerate}[label=\textbf{\alph*)}, leftmargin=*]
  \item Ordnen Sie jede Situation einer \textbf{Marktform} (Polypol, Oligopol, Monopol) zu
        und begründen Sie über \textbf{Marktmacht} und \textbf{Wettbewerb}:
  \begin{enumerate}[label=(\arabic*), leftmargin=*]
    \item viele kleine Bäckereien in einer Großstadt
    \item drei große Mobilfunknetzbetreiber in einem Land
    \item der einzige regionale Wasserversorger
  \end{enumerate}
  \item Ein Lieferdienst-Marktplatz vermittelt zwischen Restaurants und Gästen. Erklären Sie
        \textbf{Netzwerkeffekt} und \textbf{Lock-in-Effekt} mit je einem Beispiel.
  \item Nennen Sie \textbf{zwei Markteintrittsbarrieren} und je eine Folge für den Wettbewerb.
\end{enumerate}
\karo{12.5cm}

% ===================== AUFGABE 7 =====================
\aufgabe{Aufgabe 7 -- BCG, Porter und Skaleneffekte}{\diff{med}{mittel} \;(12 P)}{%
VL4 -- Unternehmensstrategie}

\begin{enumerate}[label=\textbf{\alph*)}, leftmargin=*]
  \item Nennen Sie je \textbf{einen Vorteil} und \textbf{einen Nachteil} der \textbf{BCG-Matrix}.
  \item Ordnen Sie jedem Unternehmen die passende \textbf{Wettbewerbsstrategie nach Porter}
        (Preisführerschaft, Differenzierung, Fokussierung) zu und begründen Sie:
  \begin{enumerate}[label=(\arabic*), leftmargin=*]
    \item ein Lebensmittel-Discounter
    \item ein Hersteller von Premium-Smartphones
    \item ein Anbieter von Spezialsoftware nur für Zahnarztpraxen
  \end{enumerate}
  \item \textbf{Skaleneffekte:} Die Fixkosten betragen \textbf{600.000\,€/Jahr}, die variablen
        Stückkosten \textbf{8\,€}. Berechnen Sie die \textbf{Stückkosten} bei 10.000 und bei
        60.000 Stück. Erklären Sie die \textbf{Fixkostendegression} und den Bezug zur
        Preisführerschaft.\\[1mm]
        \textit{Hinweis:} $\text{Stückkosten}=\text{var.\ Stückkosten}+\dfrac{\text{Fixkosten}}{\text{Menge}}$
\end{enumerate}
\karo{10.5cm}

% ===================== AUFGABE 8 =====================
\aufgabe{Aufgabe 8 -- SWOT und Kundennutzen}{\diff{med}{mittel} \;(11 P)}{%
VL4 -- Unternehmensstrategie}

Ein Start-up bietet eine \textbf{KI-gestützte Ernährungs-App} an.

\begin{enumerate}[label=\textbf{\alph*)}, leftmargin=*]
  \item Erstellen Sie eine \textbf{SWOT-Analyse} mit je \textbf{zwei} Punkten zu Stärken,
        Schwächen, Chancen und Risiken.
  \item Leiten Sie \textbf{zwei Normstrategien} ab (z.\,B.\ Ausbauen / Aufholen / Absichern /
        Vermeiden) und benennen Sie, welche SWOT-Felder Sie kombinieren.
  \item Definieren Sie \textbf{USP} und nennen Sie einen möglichen USP der App. Ordnen Sie den
        vier Nutzenarten -- \textbf{ökonomisch, funktional, emotional, sozial} -- je ein
        Beispiel zu.
\end{enumerate}
\karo{12cm}

% ===================== AUFGABE 9 =====================
\aufgabe{Aufgabe 9 -- Produktstrategie und Produktlebenszyklus}{\diff{med}{mittel} \;(11 P)}{%
VL4 -- Unternehmensstrategie}

\begin{enumerate}[label=\textbf{\alph*)}, leftmargin=*]
  \item Ordnen Sie jede Maßnahme einer \textbf{Produktstrategie} zu (Produktvariation,
        Produktdifferenzierung, Produktbeibehaltung, Produktelimination):
  \begin{enumerate}[label=(\arabic*), leftmargin=*]
    \item Ein Smartphone-Modell wird mit besserer Kamera neu aufgelegt und ersetzt das alte.
    \item Es werden zusätzlich eine \glqq Pro\grqq- und eine \glqq Mini\grqq-Version angeboten.
    \item Ein veraltetes Modell wird aus dem Sortiment genommen.
    \item Das bewährte Ladegerät wird unverändert weitergeführt.
  \end{enumerate}
  \item Nennen Sie die \textbf{Phasen des Produktlebenszyklus} in richtiger Reihenfolge und
        \textbf{skizzieren} Sie den typischen Verlauf von Umsatz/Absatz über die Zeit
        (im Karofeld). Ordnen Sie zwei Phasen je eine sinnvolle Maßnahme zu.
  \item Nennen Sie \textbf{einen Kritikpunkt} am Modell.
\end{enumerate}
\karo{11cm}

% ===================== AUFGABE 10 =====================
\aufgabe{Aufgabe 10 -- Konsumentenverhalten und AIDA}{\diff{med}{mittel} \;(12 P)}{%
VL5 -- Marketing}

Vergleichen Sie den Kauf eines \textbf{Laptops} mit dem eines \textbf{Kaugummis}.

\begin{enumerate}[label=\textbf{\alph*)}, leftmargin=*]
  \item Erklären Sie den Unterschied zwischen \textbf{S-R-Modell} und \textbf{S-O-R-Modell}
        (Stichwort \glqq Black Box\grqq).
  \item Ordnen Sie Laptop und Kaugummi je \textbf{High-} oder \textbf{Low-Involvement} zu und
        nennen Sie die passende \textbf{Kaufentscheidungsart} (extensiv, limitiert, impulsiv,
        habitualisiert).
  \item Nennen Sie die \textbf{drei Arten von Informationsquellen} der Informationssuche mit je
        einem Beispiel.
  \item Entwickeln Sie für eine \textbf{Marketing-Botschaft} des Laptops je eine Aussage für
        die vier \textbf{AIDA}-Stufen.
\end{enumerate}
\karo{11cm}

% ===================== AUFGABE 11 =====================
\aufgabe{Aufgabe 11 -- Beschaffung und optimale Bestellmenge}{\diff{hard}{schwer} \;(14 P)}{%
VL7 -- Beschaffung, SCM und Logistik}

\begin{enumerate}[label=\textbf{\alph*)}, leftmargin=*]
  \item Ordnen Sie jedes Gut der richtigen \textbf{Materialart} zu (Rohstoff, Hilfsstoff,
        Betriebsstoff) -- für einen Fahrradhersteller:
  \begin{enumerate}[label=(\arabic*), leftmargin=*]
    \item Aluminium für die Rahmen
    \item Schrauben und Klebstoff
    \item Schmieröl für die Fertigungsmaschinen
  \end{enumerate}
  \item \textbf{Make-or-Buy:} Nennen Sie zwei Entscheidungskriterien und entscheiden Sie
        beispielhaft, ob die Rahmen selbst gefertigt oder zugekauft werden sollten.
  \item \textbf{Optimale Bestellmenge:} Ein Bauteil wird mit \textbf{Jahresbedarf
        $m = 8000$ Stück} benötigt. Die \textbf{Bestellkosten} betragen $K_B = 40$\,€ pro
        Bestellung, der \textbf{Einstandspreis} $p = 50$\,€/Stück und der
        \textbf{Lagerhaltungskostensatz} $z = 8\,\%$ pro Jahr.\\[1mm]
        \textit{Formel (Andler):}\;
        $\displaystyle q_{\text{opt}}=\sqrt{\frac{2\cdot m\cdot K_B}{k_L}}$\;\; mit \;
        $k_L = p\cdot z$ \;(Lagerkosten je Stück und Jahr).\\[1mm]
        Berechnen Sie (i) die \textbf{optimale Bestellmenge}, (ii) die \textbf{Anzahl der
        Bestellungen} pro Jahr und (iii) prüfen Sie über
        $\text{Bestellkosten}=\text{Lagerkosten}$.
\end{enumerate}
\karo{9.5cm}

% ===================== AUFGABE 12 =====================
\aufgabe{Aufgabe 12 -- Personalbedarf und Personalmanagement}{\diff{med}{mittel} \;(11 P)}{%
VL8 -- Personalmanagement}

Die Logistikabteilung der \textit{NovaTech GmbH} plant das nächste Jahr. Wegen eines neuen
Lagers wird ein \textbf{Sollbestand (Bruttopersonalbedarf) von 18} Mitarbeitenden benötigt.
Aktuell sind \textbf{14} beschäftigt. Im nächsten Jahr: \textbf{2} gehen in Rente,
\textbf{1} wechselt in eine andere Abteilung, \textbf{1} kehrt aus der Elternzeit zurück und
\textbf{1} bereits unterschriebener neuer Vertrag beginnt.

\begin{enumerate}[label=\textbf{\alph*)}, leftmargin=*]
  \item Berechnen Sie den \textbf{fortgeschriebenen (zukünftigen) Personalbestand} und den
        \textbf{Netto-Personalbedarf}. Wie viele Personen müssen zusätzlich beschafft werden?
  \item Nennen Sie je einen \textbf{Vorteil} der \textbf{internen} und der \textbf{externen}
        Personalbeschaffung.
  \item Ordnen Sie drei Personalentwicklungsmaßnahmen den Formen zu
        (\textit{into / on / near / off the job}): Trainee-Programm, Job Rotation,
        berufsbegleitendes Studium.
\end{enumerate}
\karo{11.5cm}

\end{document}
